Nguồn: \href{https://nguoidothi.net.vn/nha-van-nguyen-ngoc-phan-dinh-dieu-da-song-mot-cuoc-doi-rat-dang-hoang-39640.html}{Báo Người Đô Thị}

\begin{quote}
“Anh Diệu luôn trân trọng con người để luôn có được những ứng xử, chăm lo tinh tế nhất khiến chúng ta chỉ còn biết nói rằng, sống tử tế là phải như thế đấy, chơi với nhau là phải như thế…” - nhà văn lão thành Nguyên Ngọc nói về người bạn lớn của mình, cũng là người từng đồng hành cùng ông trong suốt một nhiệm kỳ làm Đại biểu Quốc hội (khoá VI): GS. Phan Đình Diệu, nhân kỷ niệm 5 năm ngày mất (13.5.2018) và 87 năm ngày sinh (12.6.1936) của nhà toán học, nhà khoa học máy tính nổi tiếng.
\end{quote}

{\bf “Cái tự do thơ thới ấy, nó rất gần với văn nghệ sĩ”}

\emph{Trước khi có một tình bạn thân thiết với GS. Phan Đình Diệu, ông từng quan sát nhà toán học ở cự ly nào?}

Năm 1988, khi tôi đang là Tổng biên tập báo Văn Nghệ, thì có một bài báo gửi đến toà soạn, ngắn thôi, chỉ khoảng tầm hơn 200 chữ, ký tên hai ông Phan Đình Diệu và Hoàng Tụy.

Nội dung bày tỏ mối băn khoăn của hai nhà khoa học hàng đầu về khái niệm “quay hộp đen” mà Tổng Bí thư thời kỳ ấy đã sử dụng trong một bài phát biểu đăng báo (thậm chí về sau còn thành tên chuyên mục của tờ báo ấy).

Tác giả bài báo cho rằng, đây là một cách vận dụng sai thuật ngữ và là một phép liên tưởng thiếu chính xác, do chưa hiểu thấu đáo về kiến thức chuyên ngành. Bài viết tuy ngắn kỷ lục và chỉ chiếm một ô nhỏ trên báo nhưng thực sự gây chấn động, khiến báo bạn sau đó phải bỏ hẳn chuyên mục đó.

Nhưng trước đó thì chúng tôi cũng đã quen biết nhau, đặc biệt là khi cùng tham gia Quốc hội khóa VI (1976 - 1981), thuộc Ủy ban Văn hóa Giáo dục của Quốc hội.

Ở vào cái thời mà hội trường Quốc hội chưa có bấm nút, chưa có chất vấn Đại biểu Quốc hội… thì những tiếng nói phản biện đã cất lên trong bối cảnh thế nào?

Diễn đàn Quốc hội bấy giờ thì gần như không có tiếng nói phản biện, riêng Phan Đình Diệu luôn là người đầu tiên đưa ra những ý kiến trái chiều. Anh Diệu thường nói nhiều về việc cần lắng nghe ý kiến của trí thức, đặc biệt là vấn đề dân chủ.

Thường trong nhóm chúng tôi hồi đó, những ý kiến của các anh Phan Đình Diệu, Tạ Quang Bửu… bao giờ cũng là những ý kiến rất sâu sắc và đáng lắng nghe. Ba anh em chúng tôi hồi đấy cũng chơi thân nhau nên cũng bàn được với nhau nhiều việc, nhiều chuyện.

Là một nhà toán học nhưng anh Diệu có mối quan tâm rộng lắm, anh cũng đặc biệt quan tâm đến văn hóa. Hồi Thủ tướng Võ Văn Kiệt còn sống, ông từng gợi ý cho tôi là nên thành lập một viện nghiên cứu tư nhân về văn hóa, và người đầu tiên tôi nghĩ ngay tới không ai khác chính là GS. Phan Đình Diệu. Hai anh em cũng đã bàn đến việc nếu làm nghiên cứu về văn hóa thì nên làm thế nào để văn hóa gắn liền với giáo dục, làm sao để tạo ra được một cái phanh cho vấn nạn văn hoá đang xuống cấp trầm trọng… 

\emph{Đó có phải là một trong những lý do khiến GS. Phan Đình Diệu trở thành một cái tên có sức hút mạnh với giới văn nghệ sĩ?}

Lý do trước hết, theo tôi, anh Diệu là một con người rất đàng hoàng. Anh thực sự đã sống một cuộc đời rất đàng hoàng, với tất cả sự thẳng thắn và trung thực. Cái tự do thơ thới ấy, nó rất gần với giới văn nghệ sĩ.

Anh Diệu cũng là người làm thơ rất hay, có một bài thơ anh làm tặng GS. Trần Đại Nghĩa tôi rất thích: ...Tình nặng ấy chưng tình đất nước/Nghiệp đời há kể nghiệp vàng son...

Con người đó là người rất hiểu thơ và rất hiểu văn học, vì thực ra toán học nó rất gần với văn học. Tôi vẫn hay nói với các anh em trẻ, nhất là mấy anh em trẻ viết văn, là một trong những nhược điểm của các bạn là ít quan tâm đến khoa học tự nhiên quá, trong khi khoa học tự nhiên, chẳng hạn như toán học, cái sức tưởng tượng của nó ghê gớm lắm (con tôi cũng làm toán nên tôi biết), vì thế nó rất gần với văn học nghệ thuật. Tôi với anh Diệu dễ gần nhau, thân nhau chắc cũng là vì thế.


\emph{Điều gì đã làm nên sự đàng hoàng đó, theo ông?}

Tính cách Nghệ là một phần, cái nôi địa linh nhân kiệt là một phần. Phần nữa, quan trọng không kém, là thế hệ của chúng tôi đã may mắn được hưởng một nền giáo dục phải nói là rất đàng hoàng, với những ông thầy “ghê gớm”, có sức ảnh hưởng rất lớn tới học trò. Thế hệ bọn tôi được một cái điều như thế!

\emph{Dường như lâu lắm tôi mới nghe đến ba chữ tưởng như rất quen tai này: “sống đàng hoàng”. Theo ông, thế nào là “sống đàng hoàng”?}

Nó thật ra rất đơn giản, nó đến từ sự tử tế. Anh Diệu luôn trân trọng con người để luôn có được những ứng xử, chăm lo tinh tế nhất khiến chúng ta chỉ còn biết nói rằng: Sống tử tế là phải như thế đấy, chơi với nhau là phải như thế…

Anh Diệu là người sống rất tình cảm. Tôi nhớ anh từng kể với tôi, thời Cải cách ruộng đất, có lần anh thấy người ta giải một “tên địa chủ”, lại gần thì hóa ra là ông thầy mà anh ấy rất kính trọng, thế là cậu học trò ấy cứ đứng đấy khóc mãi. Tận tới sau này khi kể lại cho tôi nghe, anh còn ứa nước mắt. Cái ấn tượng đó hẳn đã hằn rất sâu trong tâm trí anh Diệu để sau này góp phần hình thành nên cái con người cương trực và tử tế nơi anh.

Anh Diệu, tôi nghĩ anh ấy là một trí thức lớn đúng nghĩa.


{\bf “Sự nhất quán ấy không dễ gì mà có được”.}

\emph{Thế nào là một “trí thức lớn đúng nghĩa”, theo ông? }

Trí thức theo tôi không chỉ là người học cao đâu, có nhiều người học rất ghê, nhưng người ta cũng chỉ có thể gọi họ là nhà bác học, chứ không phải là một bậc trí thức. Trí thức, nói một cách nôm na, còn phải là những người “giữa đường thấy sự bất bằng chẳng tha”; hay như từ của anh Cao Huy Thuần, là mấy người hay xớ rớ vào việc người khác, rằng việc đó không ảnh hưởng gì tới họ, nhưng họ vẫn thấy mình liên quan, mình cần có trách nhiệm tới cùng.

Nhiệm vụ của trí thức theo tôi là làm cho xã hội không được yên, không thể đứng yên, theo cái nghĩa tích cực của điều này.

\emph{Hiểu theo nghĩa đó, ông thấy xã hội Việt Nam đương đại trong mấy thập niên gần đây có được nhiều bậc trí thức? Và ông, có định vị mình là trí thức?}

Cũng có, ông Phan Đình Diệu là trí thức, ông Hoàng Tụy cũng là trí thức…, còn những người khác nữa. Tôi cũng là người trí thức, vì tôi toàn đi nói những cái chuyện có dính gì đến tôi đâu, mà đời tôi cũng chẳng cần cái gì khác ngoài cái đời sống bình thường như thế này.

\emph{Ở thời điểm bản lề trước và sau Đổi mới, khi báo Văn Nghệ (mà có thời gian ông làm Tổng biên tập) lần lượt đăng những truyện ngắn gây chấn động văn đàn của Nguyễn Huy Thiệp hay bút ký nổi tiếng Cái đêm hôm ấy… đêm gì của Phùng Gia Lộc…, thì tại diễn đàn Quốc hội (các khoá V, VI) hay Mặt trận Tổ quốc Việt Nam (các khoá III, IV, V, VI, VII), GS. Phan Đình Diệu cũng được biết là một tiếng nói phản biện đầy nhiệt thành và nhiều trăn trở. Ông nghĩ, đó là do “thời thế tạo anh hùng”, hay chính là ngược lại?}

Thực ra thì hồi đó bọn mình cũng có tranh thủ được một thời gian để kịp nói ra được những gì cần nói. Thế nhưng để có được một tiếng nói xuyên suốt, cả trong những thời điểm không thuận lợi, thì đấy còn là bản lĩnh và tầm nhìn.

Tôi thì thật ra cũng không hiểu được hết các bước trưởng thành của một tiếng nói đâu, vì với mỗi người, các bước ấy nó rất khác nhau. Có người đi nhanh hơn, nhìn thấy được sớm hơn và đáng kể, là cất được tiếng nói sớm hơn. Những tiếng nói ấy mạnh lắm, và đó là sức mạnh tự thân của họ.

Một tiếng nói như Phan Châu Trinh chẳng hạn, nó lạ lắm, tận giờ chúng ta vẫn chưa thể hiểu hết được ông ấy đâu, mấy mươi năm nay tôi cứ nghĩ mãi về cái ông Phan Châu Trinh này, rất khó hiểu vì sao lúc đó ông ấy lại có thể xuất sắc đến thế. Nhưng ông ấy cũng là một người vô cùng cô độc. Tôi nghĩ anh Diệu, anh ấy hiểu rất rõ những điều đó.  

Còn như tôi đây chẳng hạn, tôi nói thật là các bước của tôi nó cũng chậm lắm, cái giai đoạn mà tôi giáo điều nó cũng kéo dài lắm. Phải mãi đến tận năm 1979, khi sang Campuchia và tận nhìn tận thấy, tôi mới sực nhận ra, khi cái gọi là “cá nhân con người” bị đổ đồng, bị xóa sổ, bị thay thế bởi cái “chủ nghĩa bầy đàn”… thì nó có thể khủng khiếp đến thế nào.

Thế nên khi về tôi mới quyết làm cái Đề dẫn năm 1979 như sau này mọi người vẫn hay nhắc đến. Anh em văn nghệ sĩ hồi đấy hào hứng lắm, anh Diệu cũng biết chuyện bọn tôi làm cái Đề dẫn đó, và ủng hộ…

\emph{
Nhìn thấy sớm là một chuyện, nhưng hơn hết, còn là sự nhất quán, khi trước giờ vốn không thiếu những tiếng nói “linh hoạt”, “gió chiều nào xoay chiều nấy”…?}

Thực ra mà nói, cái phẩm chất đó của con người, cái sự nhất quán đó của con người, cũng không dễ gì có được, mà đôi khi không phải do họ muốn thế, tính thế. Có thể có lúc mình không biết mình sai, đời tôi sai nhiều lắm chứ, nhưng có điều mình sai một cách trung thực; mình cũng có lúc giáo điều, nhưng mình giáo điều một cách trung thực…

Anh Diệu thì khác, anh ấy có được cái sự nhất quán đó.

\emph{Ông nghĩ vì sao một tiếng nói trái chiều như Phan Đình Diệu vẫn có thể vang lên một cách dõng dạc tại các diễn đàn chính thống trong suốt một thời gian dài, dù không “dễ nghe” chút nào?}

Ở anh Diệu, anh ấy có một chữ “Liêm” rất rõ ràng, cái mục đích lên tiếng của anh ấy nó vô cùng trong sáng. Anh chỉ đơn giản là người dám nói lên sự thật thôi, còn nhiều người khác thì không dám. Anh nói rất chặt chẽ và đàng hoàng, lịch sự; không đả kích cá nhân, cũng không nói sau lưng hay bên lề. Một người đứng đắn, đàng hoàng, có uy tín quốc tế như thế, người ta phải nể chứ! 

Lê Quân thực hiện